\section{Classical comparisons and resource conventions}\label{sec:resources51}
\subsection{Positive realization and polyhedral extensions}
Benvenuti--Farina \cite[Theorem 2]{BF04} characterize stationary
positive realizability by a proper polyhedral cone containing the
input vector, lying in the observation cone, and invariant under the
signed state update. Their earlier survey
\cite[Section 3, Example 4 and Theorem 2]{BF03} is especially close
to the compatibility issue here: its periodic Hankel example admits
a nonnegative factorization of inner size three but no positive
realization of order three, and its minimality criterion includes a
nonnegative intertwining matrix. Thus neither the failure of positive
rank to determine dynamic order nor the need for positive
intertwining is new. The fixed-degree, unbounded-positive-order family
of \cite{BF99}, also discussed in \cite[Section 2]{BF03}, is another
direct antecedent.

Theorems~\ref{thm:section51} and~\ref{thm:local51} place that
principle in a controlled finite-horizon setting: there is one
stochastic update per letter and epoch, the accessible span may
change with the cut, and the projected physical polytope is obtained
from a section rather than from all hidden basis states. The
additional conclusions are the reachability-rank stratification,
the occupation bounds across separated cuts, and stationarization
from a uniform bound on independently chosen clocked machines in
Theorem~\ref{thm:finitegroup51}. The one-surplus six-state result and
the explicit positive-error occupation estimate are not consequences
of merely factoring one Hankel truncation. The finite-group proof
uses compact polytope geometry; it does not assume a single
nonnegative transition matrix in advance.

Yannakakis' factorization theorem \cite{Yannakakis} identifies the
extension complexity of a polytope with the nonnegative rank of its
slack matrix. The number of inequalities of a static extension is
not the number of vertices of its image, and it does not impose a
common stochastic extension of prescribed command maps.
Proposition~\ref{prop:dual51} makes this last requirement explicit
for supplied sections. The quarter-error example following it shows
that positivity on the section alone is insufficient. These are
uses of classical extension and separation theory, not new claims
about linear-programming duality. Likewise the simplex
volume identity in Lemma~\ref{lem:difference50} is the classical
difference-body formula \cite{RogersShephard}; the occupation
inequalities are its application to the actual stochastic process.

\subsection{Tensor completion and exact finite certificates}
Kahle--Kubjas--Kummer--Rosen \cite[Proposition 2.2 and Remarks 2.3,
2.5]{KKKR} describe rank-one completion using zero patterns and toric
relations, and explain why higher-order tensor relations need not be
square-free matrix cycles. Their Proposition 2.7 separates real sign
compatibility from complex completion for nonzero data. The
sign--magnitude theory here uses this classical algebraic structure
inside an interval-constrained, bounded, complete response table and
then compiles a common stochastic machine. It does not claim the
invention of toric circuit equations.

Gillis--Shitov \cite[Lemma 1 and Theorem 3]{GS} show that fixed-sign
rank-one matrix approximation in the infinity norm has a linear
feasibility formulation, while the unrestricted decision problem is
NP-complete. This is the relevant matrix precedent for separating
sign choices from magnitude constraints. The tensor certificate
bounds in Proposition~\ref{prop:complexity50} account for higher-order
relations and zero-containing intervals; they do not turn sign
search into a polynomial-time procedure. The displayed record count
prices circuit records, not the input, sign assignment, basis data or
algebraic isolating intervals, which must also be supplied.

For noisy observed cubic tensors, Nie--Tang--Zhou
\cite[Theorem 3.1]{NTZ26} prove a local Lipschitz recovery result for
a biquadratic optimization under rank and uniqueness assumptions;
Remark 3.2(i) distinguishes the existence of a sufficiently small
noise neighborhood from an explicit admissible threshold. That
recovery problem differs from a lower bound against every hidden
stochastic realization of a complete controlled experiment.
Lemma~\ref{lem:anchor51} derives its conditioning bound directly
from the prescribed signed seeds, and
Corollary~\ref{cor:robust-four51} gives a horizon-independent error
interval. It neither removes the hypotheses of the tensor recovery
theorem nor claims a general noisy tensor optimizer.

\subsection{Five resource questions}
The following conventions apply throughout. A read-only program is
not a free persistent private random selector: any such selector
retained between epochs belongs to the counted state.

\begin{center}
\small
\begin{tabular}{@{}p{.22\linewidth}p{.72\linewidth}@{}}
\toprule
Question & Resource and conclusion \\
\midrule
Atomic-row existence & $W_{N,\varepsilon}$ counts available persistent
labels. Horizon-specific real rows and the external clock are allowed;
program length, arithmetic and sampling work are not thereby bounded.\\[4pt]
Real-algebraic synthesis & Algebraic input permits finite elimination
and algebraic output in Theorem~\ref{thm:local51}. The input encoding
and elimination work are separate costs. The local supplied-witness
checker is not a general optimizer over unknown sections.\\[4pt]
Rational description & Numerator and denominator lengths are charged
when a rational implementation is considered. Exact algebraic optima
need not be attainable by rational rows; finite rational errors are
maxima of finitely many rational values.\\[4pt]
Finite-bit implementation & The appendices state their precision,
rounding and stored-bit conventions explicitly. Their bounds do not
follow by declaring an exact real row to occupy finitely many bits.\\[4pt]
Uniform computation & A label bound alone supplies neither a uniform
polynomial-time program nor a bounded-time exact sampler for arbitrary
real probabilities. No such deployment theorem is used here.\\
\bottomrule
\end{tabular}
\end{center}

The two-state sign algorithm takes an explicitly expanded response
table; the orthogonal encoding of an arbitrary binary table has
dimension $2^n+1$. The local projector formula instead has polynomial
description length in the listed clock and register sizes. These
statements concern different representations, neither of which
licenses a polynomial-time claim in a succinct horizon encoding.
The certificate bounds are deliberately coarse. Positive radical
roots specify chosen algebraic values; a short radical expression is
not an assertion about the degree of their full compositum.

The analytic papers in the General Theta series have independent
model-specific obligations: returned Fourier estimates, stopped
large deviations, canonical conditioning, nonlinear resolvents,
filter regularity and operator-domain response identities. The
finite-dimensional arguments in this article require none of those
results and do not establish them. This identifies the mathematical
interface without treating a dependency ledger as an analytic proof.
